\subsection{Examples of integrally closed domains}

PROPOSITION 10. Every principal ideal domain is integrally closed.\par

Let A be a principal ideal domain, K its field of fractions and x an element of K. There exist two relatively prime elements a, b of A such that \(x = ab^{-1}\) (Algebra, Chapter VII, §1, no. 2, Proposition 1 and Chapter VI, §1, no. 11, Proposition 9 (DIV)). If x is integral over A, it is a root of a polynomial

\(X^{n} + c_{1}X^{n-1} + \cdots + c_{n}\) of A{[}X{]}. Then \(a^{n} = b(-c_{1}a^{n-1} - \cdots - c_{n}b^{n-1})\) , which proves that \(b\) divides \(a^{n}\) . Since \(a\) and \(b\) are relatively prime, this implies that \(b\) is invertible in A (Algebra, Chapter VI, § 1, no. 12, Corollary 1 to Proposition 11 (DIV)); hence \(x \in A\) .

LEMMA 2. Let R be a ring and P a monic polynomial in R{[}X{]}. There exists a ring R' containing R such that in the polynomial ring R'{[}X{]} the polynomial P is a product of monic polynomials of degree 1.

We proceed by induction on the degree \(n\) of \(\mathbf{P}\) , the lemma being obvious for \(n = 0\) and \(n = 1\) . Suppose therefore that \(n > 1\) . Let \(\mathfrak{a}\) be the ideal of \(\mathbb{R}[\mathbf{X}]\) generated by \(\mathbf{P}\) and let \(f\) be the canonical homomorphism of \(\mathbb{R}[\mathbf{X}]\) onto \(\mathbf{B} = \mathbb{R}[\mathbf{X}] / \mathfrak{a}\) . Since \(\mathbf{P}\) is monic, for every polynomial \(\mathbf{Q} \in \mathbb{R}[\mathbf{X}]\) , \(\deg(\mathrm{PQ}) = \deg(\mathrm{P}) + \deg(\mathrm{Q})\) , whence \(\mathfrak{a} \cap \mathbb{R} = 0\) ; the restriction of \(f\) to \(\mathbb{R}\) is therefore injective. Identifying \(\mathbb{R}\) with the subring \(f(\mathbb{R})\) of \(\mathbb{B}\) by means of \(f\) and writing \(b = f(\mathbf{X})\) , we see that \(b\) is a root of \(\mathbf{P}\) in \(\mathbb{B}\) , \(\mathbf{P}\) being considered as a polynomial in \(\mathbb{B}[\mathbf{X}]\) . Then there exists a monic polynomial \(\mathbf{Q}\) in \(\mathbb{B}[\mathbf{X}]\) of degree \(n - 1\) such that \(\mathrm{P}(\mathbf{X}) = (\mathbf{X} - b)\mathrm{Q}(\mathbf{X})\) (Algebra, Chapter IV, § 1, no. 4, Proposition 5). By the induction hypothesis there exists a ring \(\mathbb{R}' \supseteq \mathbb{B}\) such that in \(\mathbb{R}'[\mathbf{X}]\) the polynomial \(\mathbf{Q}\) is a product of monic polynomials of degree 1; clearly in \(\mathbb{R}'[\mathbf{X}]\) \(\mathbf{P}\) is then a product of monic polynomials of degree 1.

PROPOSITION 11. Let A be a ring, R an A-algebra and P and Q monic polynomials in R{[}X{]}. If the coefficients of PQ are integral over A, the coefficients of P and Q are integral over A.

By a double application of Lemma 2, we see that there exists a ring \(\mathbf{R}'\) containing \(\mathbf{R}\) and families of elements \((a_i)_{1 \leqslant i \leqslant m}, (b_j)_{1 \leqslant j \leqslant n}\) of \(\mathbf{R}'\) such that in \(\mathbf{R}'[\mathbf{X}]\) \(\mathrm{P}(\mathbf{X}) = \prod_{i=1}^{m} (\mathbf{X} - a_i)\) , \(\mathbf{Q}(\mathbf{X}) = \prod_{j=1}^{n} (\mathbf{X} - b_j)\) ; the coefficients of \(\mathbf{PQ}\) belong to the integral closure \(\mathbf{A}'\) of \(\mathbf{A}\) in \(\mathbf{R}'\) and hence (no. 2, Proposition 7) the elements \(a_i\) ( \(1 \leqslant i \leqslant m\) ) and \(b_j\) ( \(1 \leqslant j \leqslant n\) ) belong to \(\mathbf{A}'\) . It follows that the coefficients of \(\mathbf{P}\) and \(\mathbf{Q}\) are integral over \(\mathbf{A}\) (no. 1, Corollary 1 to Proposition 4).

Let A be an integral domain, K its field of fractions and \(K'\) a K-algebra (not necessarily commutative). Given an element \(x \in K'\) algebraic over K, the polynomials \(P \in K[X]\) such that \(P(x) = 0\) form an ideal \(a \neq 0\) of \(K[X]\) , necessarily principal (Algebra, Chapter IV, §1, no. 5, Proposition 7). There exists a unique monic polynomial which generates a; generalizing the terminology introduced in Algebra, Chapter V, §3, no. 1, Definition 3, this monic polynomial will be called the minimal polynomial of x over K.

COROLLARY. Let A be an integral domain, K its field of fractions and x an element of a K-algebra \(K'\) (not necessarily commutative). If x is integral over A, the coefficients of the minimal polynomial P of x over K are integral over A (and they therefore belong to A if A is integrally closed).

There exists by hypothesis (no. 1, Theorem 1) a monic polynomial \(Q \in A[X]\) such that \(Q(x) = 0\) . As P divides Q in K{[}X{]}, it follows from Proposition 11 that the coefficients of P are integral over A.

Let A be a ring and R a (commutative) A-algebra; the homomorphism \(\phi: A \rightarrow R\) defining the A-algebra structure on R can be extended uniquely to a homomorphism \(A[X] \rightarrow R[X]\) of polynomial rings over A and R, leaving X invariant and hence \(R[X]\) is given a canonical A{[}X{]}-algebra structure.

PROPOSITION 12. Let A be a ring, R an A-algebra and P a polynomial in \(R[X_{1},\ldots,X_{n}]\) . For P to be integral over \(A[X_{1},\ldots,X_{n}]\) , it is necessary and sufficient that the coefficients of P be integral over A.

By considering the polynomials of \(R[X_{1},\ldots,X_{n}]\) as polynomials in \(X_{n}\) with coefficients in \(R[X_{1},\ldots,X_{n-1}]\) , we see immediately that it is reduced to proving the proposition for n=1. Then let P be a polynomial in \(R[X]\) ; it follows immediately from no. 1, Proposition 5 that, if the coefficients of P are in the integral closure B of A in R, the element P, which belongs to \(B[X]=B\otimes_{A}A[X]\) , is integral over A{[}X{]}. Conversely, suppose that P is integral over A{[}X{]} and let

\[
\mathrm{Q} (\mathrm{Y}) = \mathrm{Y} ^ {m} + \mathrm{F} _ {1} \mathrm{Y} ^ {m - 1} + \dots + \mathrm{F} _ {m}
\]

be a monic polynomial with coefficients \(F_{i} \in A[X]\) with P as a root. Let r be an integer strictly greater than all the degrees of the polynomials P and \(F_{i}\) ( \(1 \leqslant i \leqslant m\) ) and let us write

\[
\mathrm{P} _ {1} (\mathrm{X}) = \mathrm{P} (\mathrm{X}) - \mathrm{X} ^ {r}.
\]

Then \(\mathbf{P}_1\) is a root of the polynomial

\[
\mathrm{Q} _ {1} (\mathrm{Y}) = \mathrm{Q} (\mathrm{Y} + \mathrm{X} ^ {r}) = \mathrm{Y} ^ {m} + \mathrm{G} _ {1} \mathrm{Y} ^ {m - 1} + \dots + \mathrm{G} _ {m}
\]

with coefficients in A{[}X{]}; we may therefore write

\[
- \mathrm{P} _ {1} (\mathrm{P} _ {1} ^ {m - 1} + \mathrm{G} _ {1} \mathrm{P} _ {1} ^ {m - 2} + \dots + \mathrm{G} _ {m - 1}) = \mathrm{G} _ {m}.\tag{1}
\]

Now the choice of \(r\) implies that \(-\mathbf{P}_1\) is a monic polynomial of \(\mathbb{R}[\mathbf{X}]\) and so is \(\mathbf{G}_m(\mathbf{X}) = \mathbf{Q}(\mathbf{X}^r)\) , the degrees of the polynomials \(\mathrm{F}_k(\mathbf{X})\mathbf{X}^{r(m - k)}\) being all \(< rm\) for \(k \geqslant 1\) . We conclude first of all that the polynomial

\[
\mathrm{P} _ {1} ^ {m - 1} + \mathrm{G} _ {1} \mathrm{P} _ {1} ^ {m - 2} + \dots + \mathrm{G} _ {m - 1}
\]

of R{[}X{]} is also monic; moreover, as the coefficients of \(G_{m}\) belong to A, Proposition 11 shows that \(P_{1}\) has coefficients integral over A and the coefficients of P are therefore certainly integral over A.

PROPOSITION 13. Let A be a ring, R an A-algebra and A' the integral closure of A in R. Then the integral closure of A{[}X₁, \ldots, Xₙ{]} in R{[}X₁, \ldots, Xₙ{]} is equal to A'{[}X₁, \ldots, Xₙ{]}.

This follows from Proposition 12 and Definition 3 of no. 2.

COROLLARY 1. Let A be an integral domain and A' its integral closure. Then the integral closure of the polynomial ring A{[}X₁, \ldots, Xₙ{]} is A'{[}X₁, \ldots, Xₙ{]}.

Arguing by induction on n, the problem is immediately reduced to the case n = 1. Let K be the field of fractions of A, which is also that of A'; if an element P of the field of fractions K(X) of A{[}X{]} is integral over A{[}X{]}, it belongs to the polynomial ring K{[}X{]}, for the latter is a principal ideal domain (Algebra, Chapter IV, § 1, no. 5, Proposition 7) and hence integrally closed (Proposition 10); the corollary then follows from Proposition 13 applied to R = K.

COROLLARY 2. Let A be an integral domain. For the polynomial ring A{[}X \(_{1}\) , \ldots, X \(_{n}\) {]} to be integrally closed, it is necessary and sufficient that A be integrally closed.

COROLLARY 3. If K is a field, every polynomial algebra \(\mathbf{K}[\mathbf{X}_1,\ldots ,\mathbf{X}_n]\) is an integrally closed domain.

